% GENERATED FILE. Edit canonical lessons, then run npm run generate:books.
% canonical-source-hash: fnv1a64:d08ba60c43927140
% canonical-lessons: HI-C365-siddhant, HI-C365-vyavahar, HI-C365-apvad, HI-C365-tikna, HI-C365-manvana

\chapter{Theory and Practice}
\label{ch:hi-theory-and-practice}
\hlchaptermodality{365}

\begin{chapteropening}
\textbf{By the end of this chapter:} \emph{I can say whether an idea holds up in practice.}

Complete the last lesson of chapter 365: I can say whether an idea holds up in practice.
\end{chapteropening}

\section[\texorpdfstring{\hi{सिद्धांत} (siddhānt)}{सिद्धांत (siddhānt)}]{\hi{सिद्धांत} (siddhānt) — a principle, a theory}

\label{lesson:HI-C365-siddhant}



\begin{quote}
Before the new one: say the Hindi for an essay, then the Hindi for to pick out.
\end{quote}

\subsection*{You'll want to know: \hi{सिद्धांत}}
\textbf{\hi{सिद्धांत}} — \emph{siddhānt} — \textquotedblleft{}a principle, a theory\textquotedblright{}.

\begin{quote}
\hi{यह} \hi{सिद्धांत} \hi{किताब} \hi{में} \hi{अच्छा} \hi{लगता} \hi{है।}
\end{quote}

\begin{quote}
\emph{yah siddhānt kitāb meṁ acchā lagtā hai.}
\end{quote}

\begin{quote}
This theory looks good in a book.
\end{quote}

Masculine.

\begin{cousinweb}
\textbf{\hi{सिद्धांत}} (\emph{siddhānt}) is a tatsama, borrowed from Sanskrit \emph{siddhānta}, \textquotedblleft{}an established conclusion\textquotedblright{} (Monier-Williams).
\end{cousinweb}

\begin{grammarlens}[title={What you've built}]
One more word for giving reasons and opinions.
\end{grammarlens}

\subsection*{Your turn}
\begin{itemize}
\raggedright
  \item \emph{Say it:} \emph{siddhānt}
  \item \emph{Say it:} \emph{siddhānt}, once more
  \item \emph{Say it:} \emph{yah siddhānt kitāb meṁ acchā lagtā hai.} — This theory looks good in a book.
  \item \emph{Recall:} say the Hindi for an essay, then the Hindi for to pick out, then say \emph{siddhānt} again
\end{itemize}

\begin{tcolorbox}[breakable,colback=teal!4,colframe=teal!35!black,title={Before you move on}]
What does \hi{सिद्धांत} mean? (\textquotedblleft{}A principle, a theory\textquotedblright{}.) Say it once more.
\end{tcolorbox}
\section[\texorpdfstring{\hi{व्यवहार} (vyavahār)}{व्यवहार (vyavahār)}]{\hi{व्यवहार} (vyavahār) — conduct, behaviour; practice}

\label{lesson:HI-C365-vyavahar}



\begin{quote}
Before the new one: say the Hindi for to pick out, then the Hindi for a theory.
\end{quote}

\subsection*{You'll want to know: \hi{व्यवहार}}
\textbf{\hi{व्यवहार}} — \emph{vyavahār} — \textquotedblleft{}conduct, behaviour; practice\textquotedblright{}.

\begin{quote}
\hi{सिद्धांत} \hi{अच्छा} \hi{है}, \hi{लेकिन} \hi{व्यवहार} \hi{में} \hi{यह} \hi{काम} \hi{नहीं} \hi{करता।}
\end{quote}

\begin{quote}
\emph{siddhānt acchā hai, lekin vyavahār meṁ yah kām nahīṃ kartā.}
\end{quote}

\begin{quote}
The theory is good, but in practice it doesn't work.
\end{quote}

Masculine.

\begin{cousinweb}
\textbf{\hi{व्यवहार}} (\emph{vyavahār}) is a tatsama, borrowed from Sanskrit \emph{vyavahāra}, \textquotedblleft{}doing, practice, conduct\textquotedblright{} (Monier-Williams).
\end{cousinweb}

\begin{grammarlens}[title={What you've built}]
One more word for giving reasons and opinions.
\end{grammarlens}

\subsection*{Your turn}
\begin{itemize}
\raggedright
  \item \emph{Say it:} \emph{vyavahār}
  \item \emph{Say it:} \emph{vyavahār}, once more
  \item \emph{Say it:} \emph{siddhānt acchā hai, lekin vyavahār meṁ yah kām nahīṃ kartā.} — The theory is good, but in practice it doesn't work.
  \item \emph{Recall:} say the Hindi for to pick out, then the Hindi for a theory, then say \emph{vyavahār} again
\end{itemize}

\begin{tcolorbox}[breakable,colback=teal!4,colframe=teal!35!black,title={Before you move on}]
What does \hi{व्यवहार} mean? (\textquotedblleft{}Conduct, behaviour; practice\textquotedblright{}.) Say it once more.
\end{tcolorbox}
\section[\texorpdfstring{\hi{अपवाद} (apvād)}{अपवाद (apvād)}]{\hi{अपवाद} (apvād) — an exception}

\label{lesson:HI-C365-apvad}



\begin{quote}
Before the new one: say the Hindi for a theory, then the Hindi for conduct.
\end{quote}

\subsection*{You'll want to know: \hi{अपवाद}}
\textbf{\hi{अपवाद}} — \emph{apvād} — \textquotedblleft{}an exception\textquotedblright{}.

\begin{quote}
\hi{हर} \hi{नियम} \hi{का} \hi{कोई} \hi{अपवाद} \hi{होता} \hi{है।}
\end{quote}

\begin{quote}
\emph{har niyam kā koī apvād hotā hai.}
\end{quote}

\begin{quote}
Every rule has some exception.
\end{quote}

Masculine.

\begin{cousinweb}
\textbf{\hi{अपवाद}} (\emph{apvād}) is a tatsama, borrowed from Sanskrit \emph{apavāda}, \textquotedblleft{}a refutation; an exception\textquotedblright{} (Monier-Williams).
\end{cousinweb}

\begin{grammarlens}[title={What you've built}]
One more word for giving reasons and opinions.
\end{grammarlens}

\subsection*{Your turn}
\begin{itemize}
\raggedright
  \item \emph{Say it:} \emph{apvād}
  \item \emph{Say it:} \emph{apvād}, once more
  \item \emph{Say it:} \emph{har niyam kā koī apvād hotā hai.} — Every rule has some exception.
  \item \emph{Recall:} say the Hindi for a theory, then the Hindi for conduct, then say \emph{apvād} again
\end{itemize}

\begin{tcolorbox}[breakable,colback=teal!4,colframe=teal!35!black,title={Before you move on}]
What does \hi{अपवाद} mean? (\textquotedblleft{}An exception\textquotedblright{}.) Say it once more.
\end{tcolorbox}
\section[\texorpdfstring{\hi{टिकना} (ṭiknā)}{टिकना (ṭiknā)}]{\hi{टिकना} (ṭiknā) — to hold up, to last}

\label{lesson:HI-C365-tikna}



\begin{quote}
Before the new one: say the Hindi for conduct, then the Hindi for an exception.
\end{quote}

\subsection*{You'll want to know: \hi{टिकना}}
\textbf{\hi{टिकना}} — \emph{ṭiknā} — \textquotedblleft{}to hold up, to last\textquotedblright{}.

\begin{quote}
\hi{तुम्हारा} \hi{तर्क} \hi{ज़्यादा} \hi{देर} \hi{नहीं} \hi{टिकेगा।}
\end{quote}

\begin{quote}
\emph{tumhārā tark zyādā der nahīṃ ṭikegā.}
\end{quote}

\begin{quote}
Your argument won't hold up for long.
\end{quote}

Past (it …): \hi{वह} \hi{तर्क} \hi{नहीं} \hi{टिका}, \emph{vah tark nahīṃ ṭikā}.

\begin{cousinweb}
\textbf{\hi{टिकना}} (\emph{ṭiknā}) is inherited from \emph{*ṭakk}, \textquotedblleft{}stop\textquotedblright{}, a form CDIAL reconstructs from its descendants (CDIAL 5420).
\end{cousinweb}

\begin{grammarlens}[title={What you've built}]
One more word for giving reasons and opinions.
\end{grammarlens}

\subsection*{Your turn}
\begin{itemize}
\raggedright
  \item \emph{Say it:} \emph{ṭiknā}
  \item \emph{Say it:} \emph{ṭiknā}, once more
  \item \emph{Say it:} \emph{tumhārā tark zyādā der nahīṃ ṭikegā.} — Your argument won't hold up for long.
  \item \emph{Recall:} say the Hindi for conduct, then the Hindi for an exception, then say \emph{ṭiknā} again
\end{itemize}

\begin{tcolorbox}[breakable,colback=teal!4,colframe=teal!35!black,title={Before you move on}]
What does \hi{टिकना} mean? (\textquotedblleft{}To hold up, to last\textquotedblright{}.) Say it once more.
\end{tcolorbox}
\section[\texorpdfstring{\hi{मनवाना} (manvānā)}{मनवाना (manvānā)}]{\hi{मनवाना} (manvānā) — to get (someone) to accept}

\label{lesson:HI-C365-manvana}



\begin{quote}
Before the new one: say the Hindi for an exception, then the Hindi for to hold up.
\end{quote}

\subsection*{You'll want to know: \hi{मनवाना}}
\textbf{\hi{मनवाना}} — \emph{manvānā} — \textquotedblleft{}to get (someone) to accept\textquotedblright{}.

\begin{quote}
\hi{अरुण} \hi{ने} \hi{अपनी} \hi{बात} \hi{मीरा} \hi{से} \hi{मनवाई।}
\end{quote}

\begin{quote}
\emph{Aruṇ ne apnī bāt Mīrā se manvāī.}
\end{quote}

\begin{quote}
Arun got Mira to accept his point.
\end{quote}

Past, with \hi{ने} (he or she …): \hi{उसने} \hi{अपनी} \hi{बात} \hi{मनवाई}, \emph{usne apnī bāt manvāī}.

\begin{cousinweb}
\textbf{\hi{मनवाना}} (\emph{manvānā}) is the causative of \hi{मानना}, \textquotedblleft{}to accept\textquotedblright{}, which is inherited from Sanskrit \emph{manyate}, \textquotedblleft{}thinks\textquotedblright{} (CDIAL 9857): to make someone accept.
\end{cousinweb}

\begin{grammarlens}[title={What you've built}]
One more word for giving reasons and opinions.
\end{grammarlens}

\subsection*{Your turn}
\begin{itemize}
\raggedright
  \item \emph{Say it:} \emph{manvānā}
  \item \emph{Say it:} \emph{manvānā}, once more
  \item \emph{Say it:} \emph{Aruṇ ne apnī bāt Mīrā se manvāī.} — Arun got Mira to accept his point.
  \item \emph{Recall:} say the Hindi for an exception, then the Hindi for to hold up, then say \emph{manvānā} again
\end{itemize}

\begin{tcolorbox}[breakable,colback=teal!4,colframe=teal!35!black,title={Before you move on}]
What does \hi{मनवाना} mean? (\textquotedblleft{}To get (someone) to accept\textquotedblright{}.) Say it once more.
\end{tcolorbox}
